\clearpage
\section{선별 문제 풀이 I}
\label{sec:quartic-quintic-solutions-a}

\subsection*{문제 \thechapter.1: 다섯 사차군 판정}

\begin{solution}
우울형 사차식
\[
  X^4+pX^2+qX+r
\]
의 삼차분해식은
\[
  R_f(Z)=Z^3-2pZ^2+(p^2-4r)Z+q^2
\]
이다.

\begin{enumerate}[label=(\alph*)]
  \item $p=-10,q=0,r=1$이므로
  \[
    R_f=Z^3+20Z^2+96Z=Z(Z+8)(Z+12).
  \]
  완전히 분해되므로 $G=V_4$이다. 판별식은
  \[
    147456=384^2.
  \]

  \item $p=-4,q=0,r=2$이므로
  \[
    R_f=Z^3+8Z^2+8Z=Z(Z^2+8Z+8).
  \]
  $1+2$형이다. 한 근 $\alpha=\sqrt{2+\sqrt2}$가 모든 근을 생성하여 분해체 차수가 $4$이므로 $G=C_4$이다. 판별식은 $2048$이다.

  \item $p=q=0,r=-2$이므로
  \[
    R_f=Z^3+8Z=Z(Z^2+8).
  \]
  분해체는 $\Q(\sqrt[4]2,i)$이고 차수는 $8$이다. 따라서 $G=D_4$이다. 판별식은 $-2048$이다.

  \item $p=-7,q=-3,r=1$이므로
  \[
    R_f=Z^3+14Z^2+45Z+9.
  \]
  가능한 유리근 $\pm1,\pm3,\pm9$가 모두 실패하므로 기약이다. 판별식은
  \[
    33489=183^2
  \]
  이므로 $G=A_4$이다.

  \item $p=q=r=1$이므로
  \[
    R_f=Z^3-2Z^2-3Z+1.
  \]
  유리근이 없어 기약이고 판별식은 $257$로 제곱이 아니다. 따라서 $G=S_4$이다.
\end{enumerate}
\end{solution}

\subsection*{문제 \thechapter.2: 일반 사차식의 삼차항 제거}

\begin{solution}
\[
  f(X)=X^4+aX^3+bX^2+cX+d
\]
에 $X=Y-a/4$를 대입하고 전개한다. $Y^3$의 계수는 사라지고
\[
  f(Y-a/4)=Y^4+pY^2+qY+r
\]
가 되며
\[
  p=b-\frac{3a^2}{8},
\]
\[
  q=c-\frac{ab}{2}+\frac{a^3}{8},
\]
\[
  r=d-\frac{ac}{4}+\frac{a^2b}{16}-\frac{3a^4}{256}.
\]
따라서 삼차분해식은
\[
  Z^3-2pZ^2+(p^2-4r)Z+q^2.
\]
이 계산은 $2$가 가역인 체에서 유효하다. 원래 다항식과 평행이동된 다항식은 근들의 차이가 같으므로 판별식과 분해체가 변하지 않는다.
\end{solution}

\subsection*{문제 \thechapter.4: $C_4$와 $D_4$의 직접 구분}

\begin{solution}
(a) $f=X^4-4X^2+2$라 하자. $u=X^2$로 두면
\[
  u^2-4u+2=0
\]
이므로
\[
  u=2\pm\sqrt2.
\]
근은
\[
  \pm\alpha,\quad\pm\beta,
  \qquad
  \alpha=\sqrt{2+\sqrt2},\quad
  \beta=\sqrt{2-\sqrt2}.
\]
그런데
\[
  \sqrt2=\alpha^2-2,
  \qquad
  \alpha\beta=\sqrt2,
  \qquad
  \beta=\frac{\sqrt2}{\alpha}.
\]
따라서 모든 근이 $\Q(\alpha)$에 들어간다. $f$는 소수 $2$에 대한 아이젠슈타인 판정법으로 기약이므로
\[
  [\Q(\alpha):\Q]=4.
\]
분해체의 갈루아군은 위수 $4$이고, 판별식이 제곱이 아니므로 $V_4$가 아니다. 따라서 $C_4$이다.

(b) $f=X^4-2$의 근은 $\pm a,\pm ia$, $a=\sqrt[4]2$이다. 분해체는
\[
  L=\Q(a,i).
\]
아이젠슈타인 판정법에서 $[\Q(a):\Q]=4$이고 $\Q(a)\subset\R$이므로 $i\notin\Q(a)$. 따라서
\[
  [L:\Q]=8
\]
이고 갈루아군은 $D_4$이다.
\end{solution}

\subsection*{문제 \thechapter.5: 오차식의 순환형 사전}

\begin{solution}
인수차수는 그대로 프로베니우스 순환의 길이가 된다.
\[
\begin{array}{c|c|c|c}
\text{인수차수} & \text{순환형} & \text{위수} & \text{부호}\\
\hline
5&(5)&5&+1\\
4+1&(4)(1)&4&-1\\
3+2&(3)(2)&6&-1\\
3+1+1&(3)(1)(1)&3&+1\\
2+2+1&(2)(2)(1)&2&+1\\
2+1+1+1&(2)(1)(1)(1)&2&-1
\end{array}
\]
길이 $m$의 순환은 부호 $(-1)^{m-1}$을 가지며, 서로소인 순환들의 위수는 길이들의 최소공배수이다.
\end{solution}

\subsection*{문제 \thechapter.7: $X^5-4X+2$}

\begin{solution}
소수 $2$에 대한 아이젠슈타인 판정법으로 $f=X^5-4X+2$는 기약이다. 따라서 갈루아군은 다섯 근 위에서 추이적이고 $5$-순환을 포함한다.

(a) $f'=5X^4-4$의 실수 영점은 $\pm a$, $a=(4/5)^{1/4}$이다. 함수는
\[
  (-\infty,-a),\quad(-a,a),\quad(a,\infty)
\]
에서 각각 증가, 감소, 증가한다. 또한
\[
  f(-a)=2+\frac{16a}{5}>0,
  \qquad
  f(a)=2-\frac{16a}{5}<0.
\]
따라서 정확히 세 실근과 한 쌍의 비실 켤레근이 있다. 복소켤레는 근들 위에서 전치이다. $5$-순환과 전치를 포함하는 추이적 부분군은 $S_5$이므로 $G=S_5$이다.

(b) 법 $7$에서 주어진 두 인수의 차수는 $2,3$이고 각각 기약이다. 판별식 $-212144$는 $7$로 나누어지지 않으므로 좋은 소수이다. 따라서 $G$는 $3+2$형 원소를 포함하고, 추이적 다섯 후보 가운데 가능한 군은 $S_5$뿐이다.
\end{solution}

\subsection*{문제 \thechapter.8: $D_5$ 예제}

\begin{solution}
(a) 평행이동하면
\[
  f(Y-2)=Y^5-10Y^4+40Y^3-80Y^2+75Y-10.
\]
최고차항을 제외한 모든 계수는 $5$로 나누어지고 상수항은 $25$로 나누어지지 않으므로 아이젠슈타인 판정법에서 기약이다.

(b) $D_5$-분해식이 분리이고 $Y^2-25$를 인수로 가지므로 유리근 $\pm5$가 있다. 부분군 분해식 판정에서
\[
  G\subseteq D_5
\]
의 켤레이다. 기약성 때문에 후보는 $C_5,D_5$이다.

(c) 법 $3$ 인수분해형은 $1+2+2$이다. 두 이차식의 판별식은 모두 $2\in\F_3$이고 이는 제곱이 아니므로 두 인수는 기약이다. 판별식 $64000000$은 $3$으로 나누어지지 않아 좋은 소수이며, $G$는 $2+2+1$형 원소를 포함한다. $C_5$에는 이런 원소가 없으므로
\[
  G=D_5.
\]
\end{solution}

\subsection*{문제 \thechapter.9: $A_5$ 예제}

\begin{solution}
법 $2$에서
\[
  \bar f=X^5+X^4+X^3+X+1.
\]
$0,1$은 근이 아니고, 유일한 기약 이차식 $X^2+X+1$로 나눈 나머지는 $X+1$이다. 차수 $5$의 가약다항식이 일차인수를 갖지 않으면 반드시 $2+3$형이므로 기약 이차인수를 가져야 한다. 따라서 $\bar f$는 기약이고 $f$도 $\Q$ 위에서 기약이다.

판별식이 $207^2$이므로
\[
  G\subseteq A_5.
\]
법 $11$에서 두 일차인수와 기약 삼차인수로 분해되므로 $G$는 삼순환을 포함한다. $A_5$ 안의 추이적 후보 $C_5,D_5,A_5$ 중 삼순환을 포함하는 것은 $A_5$뿐이다. 따라서
\[
  G=A_5.
\]
\end{solution}

\subsection*{문제 \thechapter.11: $S_4$에서 $S_3$으로}

\begin{solution}
세 쌍분할을
\[
  P_1=12|34,\qquad P_2=13|24,\qquad P_3=14|23
\]
라 쓰자. 지표 순열은 쌍분할을 다시 쌍분할로 보내므로 군준동형
\[
  \rho:S_4\to S_3
\]
를 준다. $(12)$는 $P_1$을 고정하고 $P_2,P_3$을 맞바꾸며, $(123)$은 세 $P_i$를 순환시킨다. 따라서 상에는 전치와 삼순환이 모두 있고 $\rho$는 전사이다.

항등원과 세 이중전치
\[
  1,(12)(34),(13)(24),(14)(23)
\]
는 각 $P_i$를 고정하므로 $V_4\subseteq\ker\rho$이다. 전사성에서
\[
  |\ker\rho|=24/6=4
\]
이므로
\[
  \ker\rho=V_4.
\]
\end{solution}
